\documentclass{situationist-preprint}

\title{Singer's Pond and the Limits of Analogy}
\msid{ETH-06-b}
\msfamily{ETH-06}
\msdiscipline{Normative and Applied Ethics}
\mstype{Companion manuscript}
\mskeywords{Singer; drowning child; beneficence; rescue; demandingness; iteration; collective action}

\begin{document}
\maketitle

\begin{abstract}
Singer's drowning child establishes that the affluent must help the distant needy. It is less often noticed that the analogy, pressed further, establishes nothing about how much. We identify three features of the pond case, the rescuer's unique position, the singularity of the emergency, and the absence of other agents, that the case of distant need lacks, and show that each is used when the pond is iterated into a requirement to give to the margin. Remove them and the case supports a duty to do one's share of a collective task, which is the duty ETH-06-a defends, fixed at ten per cent of income.
\end{abstract}

\section{The analogy and its use}

\citet{Singer1972} argued from a principle, that if we can prevent something very bad without sacrificing anything of comparable moral importance we ought to do it, and from a case, the child drowning in a shallow pond, in which the principle's verdict is beyond dispute. The case does two jobs. It shows that there is a duty to help, and, by its implicit claim that ruined clothes are trivial beside a child's life, it suggests that the duty's limit is set only by comparable importance. Iterated, the principle yields the marginal requirement of Definition~2.2 of ETH-06-a, a share of at least half of income for any agent whose income is at least twice the level at which comparable sacrifice begins (Proposition~3.4 of ETH-06-a). ETH-06-a argues that the duty of beneficence is instead a duty to contribute a fixed share to the collective task of meeting basic need (Principle~P1), and that the share is ten per cent (Theorem~3.5). Here we ask why the analogy does not carry Singer's argument further than that.

\section{Three disanalogies}

\paragraph{Unique position.} At the pond, the passer-by is the only one who can save the child, and if he does not, the child dies. A donor to famine relief is one of hundreds of millions who can contribute, and whether any particular child is saved depends on the contributions of many. \citet{Singer1972} anticipated this and replied that the number of others who could help does not reduce one's obligation. That is correct about the existence of the duty, but the question here is its extent. When many can help, the burden of helping can be shared, and a principle for sharing it is needed; when one alone can help, it cannot be shared and none is needed.

\paragraph{Singularity.} The pond is an emergency, met once. The need of the distant poor is chronic. The verdict at the pond is elicited for a single sacrifice, and it is iteration, not the case, that yields the marginal requirement. \citet{Unger1996} constructs variants designed to show that our different responses to pond-like and distant cases are illusions of salience. Our point is different: the cases differ in structure, since one sacrifice is called for in one and an unending series in the other, and the intuition from the first does not transfer to the series. Imagine a passer-by who lives beside a pond into which a child falls every hour. That he must save the first child is evident; that he must devote his whole life to the pond is not, and few who share the original intuition share this one.

\paragraph{Absence of other agents.} At the pond, no one else is failing to help. In the case of global need, the reason so much is left to each complier is that most others give little or nothing. The marginal requirement makes the complier's burden grow with the non-compliance of others, which is the feature of Singer's view that \citet{Murphy2000} found objectionable. The pond, which contains no non-compliers, cannot tell us whether the burden should grow in this way.

\begin{proposition}\label{p:limits}
Each step from the verdict at the pond to the marginal requirement depends on one of the three features, and none of them holds of distant need.
\end{proposition}

\begin{proof}
The step from \enquote{I must help} to \enquote{I must help, regardless of what others do} uses the absence of other agents. The step from \enquote{I must make this sacrifice} to \enquote{I must make every sacrifice short of comparable importance} uses singularity, since it treats a series of sacrifices as if each were the only one. The step from \enquote{someone must bear this cost} to \enquote{I must bear it} uses unique position. Distant need involves many potential helpers, many non-compliers, and a continuing series of demands.
\end{proof}

Removing the three features leaves a case in which many agents can contribute to meeting a continuing need, some will not, and each has a duty to help. That is the structure of Principles~P1 and~P2 of ETH-06-a: a share of a collective task, fixed in advance, and large enough to meet the need if only a fraction comply. The analogy supports the existence of the duty and is silent on its extent, which must be settled by a principle for sharing a collective burden.

\begin{objection}
Then a passer-by who has given his tenth may let the child drown.
\end{objection}

\begin{reply}
No, as ETH-06-a argues. The pond retains its three features, and so retains its verdict: there the passer-by is under a duty to rescue, which ETH-18 holds to extend up to moderate cost, and which does not draw on the share owed to the collective task. Our claim is only that the pond tells us about rescue, and that its features are what make it a case of rescue.
\end{reply}

\section{Conclusion}

The drowning child shows that the affluent must help those in need. It cannot show that they must give to the margin, because the steps that would take it there use precisely the features that distinguish rescue from beneficence. What remains is a duty to do one's share, which ETH-06-a fixes at ten per cent of income.

\begin{thebibliography}{9}\small
\bibitem[Murphy(2000)]{Murphy2000} Murphy, L.~B. (2000). \emph{Moral Demands in Nonideal Theory}. Oxford University Press.
\bibitem[Singer(1972)]{Singer1972} Singer, P. (1972). Famine, affluence, and morality. \emph{Philosophy \& Public Affairs}, 1(3), 229--243.
\bibitem[Unger(1996)]{Unger1996} Unger, P. (1996). \emph{Living High and Letting Die: Our Illusion of Innocence}. Oxford University Press.
\end{thebibliography}

\end{document}
